\begin{ledgroupsized}[r]{120mm}
\footnotesize
\pstart
\noindent\textbf{Überlieferung:}
\pend
\end{ledgroupsized}
\begin{ledgroupsized}[r]{114mm}
\footnotesize
\pstart
\parindent -6mm 
\makebox[6mm][l]{\textit{E\textsuperscript{1}}}%
Erstdruck:
\cite{01023}\textit{AE} November 1684, S.\;514. Nachdr. \textsc{Lamarra/Palaia} 2005, S.\;994.\cite{01289} (Unsere Druckvorlage.)%
\pend
\end{ledgroupsized}%
\begin{ledgroupsized}[r]{114mm}
\footnotesize
\pstart
\parindent -6mm
\parindent -6mm\makebox[6mm][l]{\textit{E\textsuperscript{2}}}%
Abdruck (nach \textit{E\textsuperscript{1}}):
\cite{01421}\textit{Opuscula omnia Actis eruditorum Lipsiensibus inserta, quae ad universam mathesin, physicam, medicinam, anatomiam, chirurgiam et philologiam pertinent}, Bd. I, Venedig 1740, S.\;281; Abbildungen am Ende des Bandes.
\pend
\end{ledgroupsized}%
\begin{ledgroupsized}[r]{114mm}
\footnotesize
\pstart
\parindent -6mm\makebox[6mm][l]{\textit{E\textsuperscript{3}}}%
Abdruck (nach \textit{E\textsuperscript{1}}):
\cite{02005}\textsc{Jacob Bernoulli}, \textit{Opera},
hrsg. von G.~\textsc{Cramer}, Genf 1744,
Bd.~I, S.\;247;
Abbildungen (Tab.~VIII, Fig.~2\textendash3) zwischen S.\;290 und S.\;291.
\pend%
\end{ledgroupsized}
\selectlanguage{latin}
\vspace{8mm}
\pstart%
\noindent%
\lbrack S.\;514\rbrack\  %
\edtext{Objectio Viri,\indexN{\textls{Vanni}, Giovanni Francesco 1638\textendash1709}
ut apparet, peringeniosi, contra receptum, et mea sententia demonstratum 
Staticorum\indexS{Statici} Theorema,\indexS{Theorema Staticorum receptum et demonstratum}}{\lemma{Objectio \lbrack...\rbrack\ Theorema}%
\Cfootnote{\textsc{G.\,F. Vanni},\indexN{\textls{Vanni}, Giovanni Francesco 1638\textendash1709} 
\cite{01995}\textit{Specimen libri de momentis gravium etc. Autore I.\,F.\,V. Lucensi}, \cite{01023}\textit{AE} November 1684, S.\;511\textendash513.}} 
non contemnenda est
quidem, solvi tamen omnino potest negando momenta in planis\indexS{momentum in plano} \textit{XC} et \textit{ZC}
in unum posse addi,\indexS{additio momentorum} ut componant momentum
gravis absolutum;\indexS{momentum gravis absolutum} 
et fraudi Viro docto\indexN{\textls{Vanni}, Giovanni Francesco 1638\textendash1709} 
fuisse videtur, quod illa vocavit partialia\indexS{momentum partiale}, hoc totale.\indexS{momentum totale}
Quid enim si grave\indexS{grave a duobus planis sustentatum} sustentetur a duobus planis
\textit{XC} inclinato,\indexS{planum inclinatum} et \textit{AC} verticali?\indexS{planum verticale} Utique momenta
in ambobus planis\indexS{momentum in plano} in unum addita\indexS{additio momentorum} non possunt
aequari uni ex ipsismet, totum parti; quod tamen
secundum objicientis\indexN{\textls{Vanni}, Giovanni Francesco 1638\textendash1709} sententiam fieri deberet: momentum enim in plano
verticali\indexS{momentum in plano verticali} 
utique est ipsum momentum gravis absolutum.\indexS{momentum gravis absolutum}
\pend
\begin{edfigure}[1.5em]%
\hfill
\includegraphics[scale=1]{%
gesamttex/edit_VIII-4/images/AE_1684_514_d1.pdf}
\hspace*{15mm}
\raisebox{0.6em}{%
\includegraphics[scale=1]{%
gesamttex/edit_VIII-4/images/AE_1684_514_d2.pdf}}
\hfill
\edcaption{%
\pagelabel{Fig:AE_1684_514_fig-1-2}%
\edtext{\lbrack\textit{Fig.~1, erg. Hrsg. nach Vanni\indexN{\textls{Vanni}, Giovanni Francesco 1638\textendash1709}}\rbrack}{%
\lemma{\lbrack\textit{Fig.~1}\rbrack}%
\Cfootnote[nonum]{%
\cite{01995}\textsc{Vanni},\indexN{\textls{Vanni}, Giovanni Francesco 1638\textendash1709} a.\,a.\,O., S.\;512.}}%
\hspace*{20mm}%
\lbrack\textit{Fig.~2}\rbrack\hspace*{10mm}}
\end{edfigure}
